\documentclass[aspectratio=169,10pt,spanish]{beamer}

\input{../../../_preambulo_beamer.tex}
\input{../../../_comandos_beamer.tex}

\selectlanguage{spanish}

\title{MA1001B - Modelación Estadística para la Toma de Decisiones}
\subtitle{Lección 39: Dos proporciones en pruebas A/B}
\author{}
\institute{Tecnol\'ogico de Monterrey}
\date{}

\begin{document}

\begin{frame}[plain]
  \vspace{1.2cm}
  {\Large\bfseries MA1001B - Modelación Estadística para la Toma de Decisiones\par}
  \vspace{0.7cm}
  {\large Lección 39: Dos proporciones en pruebas A/B\par}
  \vspace{0.7cm}
  {\color{TecRojo}\rule{\textwidth}{1.2pt}}
  \vspace{0.8cm}

  Facultad MA1001B\\[0.3cm]
  Periodo\\[0.3cm]
  Tecnol\'ogico de Monterrey
\end{frame}

\begin{frame}{La pregunta de conversión de dos versiones}
  \begin{alertblock}{Pregunta guía}
    ¿Cómo comparamos tasas de conversión de dos brazos A/B independientes, y
    qué falla si espiamos $p$ a mitad del experimento?
  \end{alertblock}

  \vspace{0.6em}
  Al final de esta lección, usted puede:
  \begin{itemize}
    \item calcular $\hat{p}_A$, $\hat{p}_B$ y la $\hat{p}$ combinada;
    \item formar el estadístico $z$ para $p_A-p_B$;
    \item decidir $H_0:p_A=p_B$ a partir de un p-valor bilateral;
    \item explicar por qué espiar infla el error tipo I.
  \end{itemize}

  \vfill
  \scriptsize Todos los visitantes e indicadores de conversión usados hoy son sintéticos. No se usan datos reales de empresa.
\end{frame}

\begin{frame}{Dos brazos de conversión independientes}
  \small
  Una prueba sintética de pago registra $x_A=72$ conversiones en $n_A=400$
  visitantes y $x_B=61$ conversiones en $n_B=410$ visitantes.

  \[
    \hat{p}_A=\frac{72}{400}=0.18,\qquad
    \hat{p}_B=\frac{61}{410}=0.148780,\qquad
    \hat{p}=\frac{133}{810}=0.164198.
  \]

  \begin{exampleblock}{Hipótesis bilaterales}
    $H_0:p_A=p_B$ frente a $H_1:p_A\neq p_B$, con $\alpha=0.05$. Los brazos
    son muestras independientes, no visitantes pareados.
  \end{exampleblock}
\end{frame}

\begin{frame}[fragile]{Paso 1: Tasas muestrales}
  \begin{columns}[T]
    \begin{column}{0.42\textwidth}
      \small
      \begin{lstlisting}
phat_a = 72 / 400
phat_b = 61 / 410
p_pool = (72 + 61) / 810
      \end{lstlisting}

      \begin{block}{Salida verificada}
        $\hat{p}_A=0.180000$\\
        $\hat{p}_B=0.148780$\\
        $p_{\mathrm{comb}}=0.164198$
      \end{block}
    \end{column}
    \begin{column}{0.58\textwidth}
      \centering
      \includegraphics[width=\textwidth]{../figuras/dos_proporciones_01_tasas_muestrales.png}
      \vspace{0.2em}
      \scriptsize Tasas de conversión observadas del Paso 1
    \end{column}
  \end{columns}
\end{frame}

\begin{frame}{$z$ combinada para $p_A-p_B$}
  \small
  \[
    z=\frac{\hat{p}_A-\hat{p}_B}{\sqrt{\hat{p}(1-\hat{p})(1/n_A+1/n_B)}}
    =\frac{0.031220}{0.026035}=1.199141.
  \]
  El p-valor bilateral es $0.230473$. Como $|z|<z^*=1.959964$ y
  $p>\alpha$, la prueba no rechaza $H_0:p_A=p_B$.
\end{frame}

\begin{frame}[fragile]{Paso 2: Decisión de la $z$ combinada}
  \begin{columns}[T]
    \begin{column}{0.40\textwidth}
      \small
      \begin{lstlisting}
se = sqrt(p*(1-p)*(1/n_a + 1/n_b))
z = (phat_a - phat_b) / se
p = 2 * norm.sf(abs(z))
      \end{lstlisting}

      \begin{block}{Salida verificada}
        $EE=0.026035$, $z=1.199141$\\
        $p=0.230473$\\
        Decisión: no rechazar $H_0$
      \end{block}
    \end{column}
    \begin{column}{0.60\textwidth}
      \centering
      \includegraphics[width=\textwidth]{../figuras/dos_proporciones_02_z_combinada.png}
      \vspace{0.2em}
      \scriptsize Colas $z$ bilaterales del Paso 2
    \end{column}
  \end{columns}
\end{frame}

\begin{frame}{Paso 3: Espiar infla el error tipo I}
  \begin{columns}[T]
    \begin{column}{0.42\textwidth}
      \small
      Simulación con semilla 42 de $8000$ pruebas A/B nulas, cada una con
      probabilidad de conversión compartida $0.16$.

      \vspace{0.4em}
      Una mirada final: tasa tipo I $=0.047250$.

      \vspace{0.3em}
      Cuatro miradas (25, 50, 75, 100 por ciento): tasa tipo I $=0.124250$.

      \vspace{0.5em}
      \textbf{Límite:} detenerse en el primer $p<0.05$ no es una prueba
      con $\alpha=0.05$.
    \end{column}
    \begin{column}{0.58\textwidth}
      \centering
      \includegraphics[width=\textwidth]{../figuras/dos_proporciones_03_limite_espiar.png}
      \vspace{0.2em}
      \scriptsize Inflación tipo I al espiar, del Paso 3
    \end{column}
  \end{columns}
\end{frame}

\begin{frame}{Verificación de conceptos}
  \begin{enumerate}
    \item Calcule $\hat{p}_A$ a partir de $x_A=72$ y $n_A=400$.
    \item ¿Por qué el EE de la prueba usa una $\hat{p}$ combinada y no dos
      EE de Wald separados?
    \item Si $p=0.230473$ y $\alpha=0.05$, ¿cuál es la decisión A/B?
    \item ¿Por qué cuatro miradas elevan la tasa de falsos positivos a
      $0.124250$?
  \end{enumerate}

  \vspace{0.6em}
  \begin{block}{Conexión con la Lección 40}
    La sesión puede continuar directamente con una prueba de bondad de
    ajuste ji-cuadrada para más de dos categorías.
  \end{block}

  \vfill
  \scriptsize Fuente: OpenStax, \textit{Introductory Business Statistics 2e},
  Capítulo 10, Sección 10.4. CC BY-NC-SA.
\end{frame}

\appendix



\end{document}
